\section{Nonaffine primitive components and compressed quotients}
\label{sec:primitive-compression}

\begin{lemma}[Core-free subgroups of a simple power]
\label{lem:simple-power-index}
Let $S$ be nonabelian simple, let $\ell$ be its least proper subgroup
index, and let $H\leq S^a$ be core-free. Then $[S^a:H]\geq\ell^a$.
\end{lemma}
\begin{proof}
We use $\ell^2<|S|$, proved in
\cite[Theorem~A(4)]{BallesterBolinchesEstebanRomeroJimenezSeral2022}.
Induct on $a$. Let $P\leq S^{a-1}$ and $B\leq S$ be the projections
of $H$, with common Goursat quotient $Q$. If $B<S$, then $|Q|<|S|$.
A full coordinate factor contained in $P$ must map trivially to $Q$,
and would therefore be contained in $H$, contrary to core-freeness.
Thus $P$ is core-free and the induction hypothesis gives
$[S^a:H]\geq[S^{a-1}:P][S:B]\geq\ell^a$.

If $B=S$, then $Q=S$: otherwise the last factor is contained in $H$.
Every full coordinate factor contained in $P$ maps onto $Q$, so there
is at most one, since images of distinct factors commute. If there is
none, the index is $|S|[S^{a-1}:P]\geq\ell^a$. If there is one,
write $P=S_i\times P'$; the map to $Q$ is an isomorphism on $S_i$
and trivial on $P'$, and $P'$ is core-free in $S^{a-2}$. Hence
$[S^a:H]=|S|[S^{a-2}:P']>\ell^a$. This proves the induction.
\end{proof}

\begin{theorem}[Primitive compression]\label{thm:primitive-compression}
If $L\leq S_r$ is nonaffine primitive, its socle $E$ is semisimple
and $L/E$ has a specified faithful action of positive degree $u$
with $2u\leq h(r)$. This assertion uses the finite primitive
classification only in actual degrees $5\leq r\leq29$.
\end{theorem}
\begin{proof}
Every nontrivial normal subgroup of a primitive group is transitive.
An abelian minimal normal subgroup is regular and gives the affine
case. Distinct minimal normal subgroups commute; since the centralizer
of a transitive group is semiregular, two such subgroups must be regular
mutual centralizers. They are isomorphic, and a third is impossible.
Consequently $E=S^a$ for a nonabelian simple $S$. Also $C_L(E)=1$:
a nontrivial centralizer would contain a minimal normal subgroup of
$L$, whereas its intersection with the centreless $E$ is trivial.
Thus
\[
 L/E\leq\operatorname{Out}(S)\wr S_a.
\]
Put $q=|\operatorname{Out}(S)|$. Its regular action gives the wreath
product a faithful action of degree $aq$, including $q=1$.

Theorem~A(4) of the cited paper gives $q\leq3\log\ell$, where
$\ell\geq5$ is the least proper subgroup index of $S$.
Lemma~\ref{lem:simple-power-index} gives $r\geq\ell^a$.
For $a\geq2$, $2aq\leq\ell^a$: at $a=2,\ell=5$, integrality
gives $q\leq6$ and $4q\leq24<25$; for $\ell\geq6$ use
$12\log\ell\leq\ell^2$. Increasing $a$ multiplies the right side
by at least five and the left side by at most $3/2$.
For $a=1,\ell\geq30$, use $6\log\ell\leq\ell$ instead.

In the remaining almost simple case, failure of the regular action
of the actual quotient to have $2|L/E|\leq r$ would imply
$r<2|L/E|\leq6\log\ell<30$. The finite primitive certificate
therefore checks precisely the required bounded range. Its complete
degree-$5$--$29$ primitive alphabet has 116 simple nonabelian socles;
each carries the literal socle and an actual regular quotient map
with $2|L/E|\leq r$. Completeness is relative to the stated primitive
classification \cite{RoneyDougal2005Primitive,PrimGrp344}; the four
nonsimple nonaffine socles in this range are covered by the preceding
argument. Thus $2u\leq r$, and parity gives $2u\leq h(r)$.
\end{proof}

\begin{theorem}[Complete semisimple compression]
\label{thm:semisimple-compression}
Suppose an actual transitive $U\leq S_w$, $w\geq5$, has a normal
semisimple subgroup $E$, possibly trivial, such that $U/E$ has a
faithful action of degree $v_0\leq h(w)/2$. The complete family of
these actions satisfies Theorem~\ref{thm:parametric}, with
$\eta_U=0$ and $\rho=1/20$.
Every action with a nonaffine primitive exact local component has
such a certificate.
\end{theorem}
\begin{proof}
Write $E=\prod_iS_i$. Each $S_i$ is a literal subgroup of $S_w$.
A nontrivial orbit, followed by a maximal overgroup of its point
stabilizer, supplies a faithful primitive action of $S_i$ of degree
at most $w$. Hence $d(S_i)\leq\log w$ by
\cite[Theorem~1.1]{HoltRoneyDougal2013}, and
$|\operatorname{Aut}(S_i)|\leq|S_i|^{\log w}$.
Since $\sum_i\log|S_i|\leq\log(w!)\leq w\log w$, the product in
Lemma~\ref{lem:semisimple-fibre} has logarithm at most
$Cw\log^2(w+b+1)$ for an absolute $C$. The reserve is
$(h(w)-v_0)/8\geq h(w)/16\geq w/20$ for $w\geq5$.
Theorem~\ref{thm:parametric} applies with $\beta=0$.
Here Lemma~\ref{lem:semisimple-fibre} is applied to the complete sum over
all literal normal axes before the physical transfer.  Its single
outer-factor product is therefore paid once.  Assigning that complete bound
to every axis and summing again would introduce a spurious normal-subgroup
multiplicity and is not part of the estimate.

For a block system of size $r$ and count $s$, let $L$ be its exact
primitive component and put $E_0=\operatorname{Soc}(L)$,
$E=U\cap E_0^s$. Each coordinate projection of $E$ is normal in
the exact $L$, and hence is a normal subproduct of $E_0$. The
subdirect-product description of products of nonabelian simple groups
makes $E$ semisimple, allowing proper diagonal and zero intersections.
Moreover $U/E$ embeds into $(L/E_0)\wr T$, whose faithful action
has degree $v_0=us$. Theorem~\ref{thm:primitive-compression} gives
$2v_0\leq h(r)s\leq h(rs)$. This proves the last assertion,
including primitive actions by taking $s=1$.
\end{proof}
